% ID: SEL-D0025-OA-E04
% ESTADO: cotejada. Segunda pasada visual completa y borrador compilado revisados.
% Fuente: Selectividad/matematicas_aplicadas/suplemente_junio_EXAMEN 4-ANDALUCÍA.pdf
% SHA-256: c1a3a7a37d38b04359115c64762b1f8a543e4d66155338931de84ef6c304da2c
% Página: 1; opción A; ejercicio original 4.

Se desea estimar la proporción de jóvenes que ven una serie de televisión. Para ello, se toma una muestra aleatoria de 100 jóvenes, de los que 36 ven la serie.
\begin{enumerate}[label=\alph*)]
\item \textbf{(1.5 puntos)} Determine un intervalo de confianza, al 96\%, para la proporción de jóvenes que ven la serie.
\item \textbf{(1 punto)} Con el mismo nivel de confianza, si queremos que el error máximo sea inferior a 0.03, \textquestiondown{}qué tamaño muestral mínimo debemos tomar?
\end{enumerate}
